\section{Standardizace dat}\label{sec:ai-standardizace}

Různé statistické znaky mohou používat odlišné jednotky a~měřítka. Věk člověka se může pohybovat v~desítkách let, zatímco jeho roční příjem ve statisících korun. Metoda založená na vzdálenosti nebo velikosti koeficientů pak může bez úpravy dat přikládat příjmu nepřiměřeně velký význam.

\begin{definition}\label{def:ai-standardizace}
    Nechť \(\mathbf{X}=(x_1,\ldots,x_n)\) a~\(\sigma_{\mathbf{X}}>0\).
    \emph{Standardizovanou hodnotou} \(x_i\) rozumíme číslo
    \[
        z_i
        =
        \frac{x_i-\average{\mathbf{X}}}{\sigma_{\mathbf{X}}}.
    \]
    Standardizovaný statistický soubor značíme
    \[
        \mathbf{Z}=(z_1,\ldots,z_n).
    \]
\end{definition}

Hodnota \(z_i\) udává, o~kolik směrodatných odchylek leží \(x_i\) nad nebo pod průměrem. Kladné hodnoty leží nad průměrem, záporné pod ním a~hodnota \(0\) odpovídá přímo průměru.

\begin{figure}[ht]
    \centering
    \begin{tikzpicture}[x=1cm,y=1cm,every node/.style={font=\normalsize}]
    \draw[diredge] (0,1.8) -- (10,1.8);
    \node[left] at (0,1.8) {\(\mathbf{X}\)};
    \foreach \x/\lab in {1/\(x_1\),3/\(x_2\),5/\(x_3\),7/\(x_4\),9/\(x_5\)}{
        \fill[blue!75!black] (\x,1.8) circle (2.3pt);
        \node[above] at (\x,1.92) {\lab};
    }
    \draw[red!75!black,very thick] (5,1.45) -- (5,2.15);
    \node[red!75!black,below] at (5,1.42) {\(\average{\mathbf{X}}\)};
    \draw[diredge] (5,1.05) -- node[right]
        {\(z_i=(x_i-\average{\mathbf{X}})/\sigma_{\mathbf{X}}\)} (5,0.35);
    \draw[diredge] (0,-0.2) -- (10,-0.2);
    \node[left] at (0,-0.2) {\(\mathbf{Z}\)};
    \foreach \x/\lab in {2/\(-1{,}4\),3.5/\(-0{,}7\),5/\(0\),6.5/\(0{,}7\),8/\(1{,}4\)}{
        \fill[green!55!black] (\x,-0.2) circle (2.3pt);
        \node[below] at (\x,-0.32) {\lab};
    }
    \draw[red!75!black,very thick] (5,-0.55) -- (5,0.15);
\end{tikzpicture}

    \caption{Standardizace posune průměr do nuly a~změní měřítko.}
    \label{fig:ai-standardizace}
\end{figure}

\begin{proposition}\label{prop:ai-standardizace}
    Pro standardizovaný soubor \(\mathbf{Z}\) platí
    \[
        \average{\mathbf{Z}}=0,
        \qquad
        \Var(\mathbf{Z})=1,
        \qquad
        \sigma_{\mathbf{Z}}=1.
    \]
\end{proposition}

\begin{proof}
    Z~\smartref[tvrzení~][]{\showrefs}{prop:ai-prumer-vlastnosti}{vlastností aritmetického průměru} dostáváme
    \[
        \average{\mathbf{Z}}
        =
        \frac1n\sum_{i=1}^{n}
        \frac{x_i-\average{\mathbf{X}}}{\sigma_{\mathbf{X}}}
        =
        \frac{1}{n\sigma_{\mathbf{X}}}
        \sum_{i=1}^{n}(x_i-\average{\mathbf{X}})
        =
        0.
    \]
    Proto
    \begin{align*}
        \Var(\mathbf{Z})
        &=
        \frac1n\sum_{i=1}^{n}z_i^2
        =
        \frac1n\sum_{i=1}^{n}
        \left(
            \frac{x_i-\average{\mathbf{X}}}{\sigma_{\mathbf{X}}}
        \right)^2\\
        &=
        \frac{\Var(\mathbf{X})}{\sigma_{\mathbf{X}}^2}
        =
        1.
    \end{align*}
    Odtud \(\sigma_{\mathbf{Z}}=\sqrt{\Var(\mathbf{Z})}=1\).
\end{proof}

\warningbox{Průměr a~směrodatnou odchylku je nutné vypočítat pouze z~trénovacích dat. Stejné dvě hodnoty potom použijeme i~pro validaci, testování a~nové objekty. Kdybychom do výpočtu zahrnuli testovací data, předali bychom modelu informaci, kterou při skutečném nasazení nemá mít.}
